\documentclass[12pt]{article}
\usepackage[utf8]{inputenc}
\usepackage[T1]{fontenc}
\usepackage[margin=1in]{geometry}
\usepackage{setspace}
\usepackage{xcolor}
\usepackage[colorlinks=true, urlcolor=blue!60!black, linkcolor=blue!60!black]{hyperref}
\setlength{\parindent}{2em}
\setlength{\parskip}{1.25em}
\renewcommand{\baselinestretch}{1.0}
\begin{document}
\begin{center}
{\LARGE \textit{Carson v. Makin}, 596 U.S. 767 (2022)\par}
\vspace{0.8em}
{\large Case Brief\par}
\end{center}
\vspace{1em}

\section{Memory Jogger}

In \href{https://www.courtlistener.com/opinion/6479298/carson-v-makin/}{\textit{Carson v. Makin}}, the Supreme Court held 6--3 that Maine's exclusion of ``sectarian'' religious schools from an otherwise generally available rural tuition-assistance program violates the First Amendment's Free Exercise Clause, rejecting any constitutional safe harbor for discrimination framed around religious ``use'' rather than religious ``status.''

\section{Detailed Case Facts}

Maine's Constitution requires the state legislature to ensure that towns provide for the support and maintenance of public schools (Me. Const., Art. VIII, pt. 1, \S 1), and state law guarantees every school-age child the opportunity for a free public education (Me. Rev. Stat. Ann., Tit. 20--A, \S 2(1)). Because Maine is the most rural State in the country, however, fewer than half of its 260 school administrative units (SAUs) operate their own public secondary schools (\href{https://www.law.cornell.edu/supremecourt/text/20-1088}{Cornell LII}).

To serve students in rural districts, Maine created a tuition-assistance program. When an SAU neither operates a public high school nor contracts with a particular public or private school to educate its residents, the SAU must pay tuition---up to a statutory cap---at the public school or approved private school selected by the parent at which the student is accepted (Me. Rev. Stat. Ann., Tit. 20--A, \S 5204(4)). Private schools qualify for approval either through accreditation by the New England Association of Schools and Colleges (NEASC) or through separate approval by the Maine Department of Education. Eligible private schools may be located out of state or abroad, may be single-sex, need not hire state-certified teachers, need not admit every applicant, and---when NEASC-accredited---are exempt from Maine's public-school curricular mandates (\href{https://www.law.cornell.edu/supremecourt/text/20-1088}{Cornell LII}).

Before 1981, parents could direct tuition payments to religious schools; during the 1979--1980 school year, more than 200 Maine students did so. In 1981, following a 1980 Maine Attorney General opinion asserting that public funding of religious schools violated the federal Establishment Clause, the Maine Legislature enacted Me. Rev. Stat. Ann., Tit. 20--A, \S 2951(2), restricting tuition assistance to ``nonsectarian'' schools. Even after the U.S. Supreme Court held in \textit{Zelman v. Simmons-Harris}, 536 U.S. 639 (2002), that neutral indirect-aid school-choice programs satisfy the Establishment Clause, the Maine Legislature declined to repeal the restriction. In administering \S 2951(2), the Maine Department of Education treated a school as ``sectarian'' if it was associated with a faith or belief system and promoted that faith or taught academic subjects through the lens of that faith (\href{https://www.law.cornell.edu/supremecourt/text/20-1088}{Cornell LII}).

Petitioners were two Maine families living in rural SAUs that neither ran a public secondary school nor contracted with a nearby school. David and Amy Carson (of Glenburn) enrolled their daughter at Bangor Christian Schools (BCS), a ministry of Bangor Baptist Church, paying tuition themselves because BCS was barred as sectarian. Troy and Angela Nelson (of Palermo) sent their son to Temple Academy (affiliated with Centerpoint Community Church) and wished to send their daughter there as well, but could not afford private tuition for both children and instead enrolled their daughter at secular Erskine Academy. Both BCS and Temple Academy held NEASC accreditation and met Maine's compulsory-attendance standards, while integrating Christian doctrine across their curricula and maintaining faith-based admissions and hiring policies. Absent \S 2951(2), both families would have directed SAU tuition funds to those schools (\href{https://www.law.cornell.edu/supremecourt/text/20-1088}{Cornell LII}).

\section{Procedural History}

\begin{itemize}
  \item \textbf{U.S. District Court for the District of Maine:} In 2018, the Carsons and Nelsons (along with the Gillis family) sued A. Pender Makin, Commissioner of the Maine Department of Education, asserting that \S 2951(2) violated the Free Exercise and Establishment Clauses of the First Amendment and the Equal Protection Clause of the Fourteenth Amendment. On cross-motions for summary judgment on a stipulated record, the District Court entered judgment for Commissioner Makin in \textit{Carson v. Makin}, 401 F. Supp. 3d 207 (D. Me. 2019), reasoning that prior First Circuit precedent upholding Maine's ban (\textit{Eulitt v. Maine Dept. of Ed.}, 386 F.3d 344 (1st Cir. 2004)) remained binding after \textit{Trinity Lutheran Church of Columbia, Inc. v. Comer}, 582 U.S. 449 (2017) (\href{https://www.law.cornell.edu/supremecourt/text/20-1088}{Cornell LII}).
  \item \textbf{U.S. Court of Appeals for the First Circuit:} While the appeal was pending, the Supreme Court decided \textit{Espinoza v. Montana Department of Revenue}, 591 U.S. 464 (2020), invalidating Montana's exclusion of church-controlled schools from a scholarship program. Although the First Circuit concluded that \textit{Espinoza} required fresh review of \textit{Eulitt}, it affirmed judgment for Maine in \textit{Carson v. Makin}, 979 F.3d 21 (1st Cir. 2020). The court distinguished \textit{Trinity Lutheran} and \textit{Espinoza} on two grounds: (1) Maine's exclusion turned on religious ``use'' of funds for sectarian instruction rather than religious ``status,'' triggering only rational-basis review; and (2) Maine's tuition benefit was meant to provide the ``rough equivalent'' of a secular public-school education (\href{https://www.law.cornell.edu/supremecourt/text/20-1088}{Cornell LII}).
  \item \textbf{U.S. Supreme Court:} The Court granted certiorari on July 2, 2021 (No. 20--1088), heard oral argument on December 8, 2021, and issued its decision on June 21, 2022 (\href{https://www.law.cornell.edu/supct/cert/20-1088}{Cornell LII Bulletin}).
  \item \textbf{Disposition:} Reversed the judgment of the First Circuit and remanded for further proceedings consistent with the opinion (\href{https://www.courtlistener.com/opinion/6479298/carson-v-makin/}{CourtListener}).
\end{itemize}

\section{Judicial Votes}

The Court ruled \textbf{6--3} in favor of the petitioners (\href{https://www.law.cornell.edu/supremecourt/text/20-1088}{Cornell LII}):

\begin{itemize}
  \item \textbf{Majority (6):} Chief Justice Roberts delivered the opinion of the Court, joined by Justices Thomas, Alito, Gorsuch, Kavanaugh, and Barrett.
  \item \textbf{Concurrences:} None.
  \item \textbf{Dissents (3):}
    \begin{itemize}
      \item Justice Breyer filed a dissenting opinion, joined in full by Justice Kagan, and joined by Justice Sotomayor as to all parts except Part I--B.
      \item Justice Sotomayor filed a separate dissenting opinion.
    \end{itemize}
\end{itemize}

\section{Holding}

\textbf{Judgment:} Maine's ``nonsectarian'' requirement for otherwise generally available secondary-school tuition-assistance payments violates the Free Exercise Clause of the First Amendment (\href{https://www.law.cornell.edu/supremecourt/text/20-1088}{Cornell LII}).

\textbf{Rule Announced:} A State has no constitutional duty to subsidize private education, but once it chooses to offer generally available public funding for private schooling, the Free Exercise Clause subjects any exclusion of religious schools---whether characterized as discrimination based on religious ``status'' or based on religious ``use'' of the funds---to strict scrutiny. A State's preference for stricter separation of church and state than the federal Establishment Clause mandates is not a compelling governmental interest (\href{https://www.law.cornell.edu/supremecourt/text/20-1088}{Cornell LII}).

\textbf{Holding vs. Dicta:} The binding holding invalidates Maine's categorical exclusion of sectarian schools from its parent-directed tuition program, rejects the status-versus-use distinction as a bypass around strict scrutiny, and confines \textit{Locke v. Davey}, 540 U.S. 712 (2004), to its narrow focus on vocational ministry degrees. Several issues lie outside the holding: (1) in footnote 1*, the majority noted that Maine did not defend its ban on the basis of BCS's or Temple Academy's specific admissions or hiring rules, leaving open whether neutral antidiscrimination conditions may govern participating schools; and (2) the Court's examples of secular public-education alternatives (such as state boarding schools, expanded busing, or remote learning) are explanatory dicta, and the Court did not resolve whether direct SAU contracts with religious schools or religious public charter schools satisfy the Establishment Clause (\href{https://www.law.cornell.edu/supremecourt/text/20-1088}{Cornell LII}; \href{https://www.cato.org/supreme-court-review/2021-2022/use-nondistinction}{Cato Supreme Court Review}).

\section{Analysis of Opinions}

\subsection{Majority Opinion (Chief Justice Roberts)}

Chief Justice Roberts treated the case as a direct application of \textit{Trinity Lutheran} and \textit{Espinoza}. Because the Free Exercise Clause bars indirect penalties and benefit exclusions as well as outright prohibitions (\textit{Lyng v. Northwest Indian Cemetery Protective Assn.}, 485 U.S. 439 (1988); \textit{Sherbert v. Verner}, 374 U.S. 398 (1963)), disqualifying otherwise eligible private schools from a generally available tuition benefit because they are religious triggers strict scrutiny (\href{https://www.law.cornell.edu/supremecourt/text/20-1088}{Cornell LII}).

\textbf{Strict Scrutiny Analysis:} Under \textit{Church of the Lukumi Babalu Aye, Inc. v. Hialeah}, 508 U.S. 520 (1993), laws targeting religious conduct survive only in rare cases where they are narrowly tailored to interests of the highest order. Because parents independently select where Maine's tuition funds go, including religious schools does not offend the federal Establishment Clause under \textit{Zelman v. Simmons-Harris}, 536 U.S. 639 (2002). Consistent with \textit{Trinity Lutheran}, \textit{Espinoza}, and \textit{Widmar v. Vincent}, 454 U.S. 263 (1981), Maine's desire to maintain stricter church-state separation than the Federal Constitution requires cannot constitute a compelling interest (\href{https://www.law.cornell.edu/supremecourt/text/20-1088}{Cornell LII}).

\textbf{Rejection of the ``Public-Education Substitute'' Characterization:} Chief Justice Roberts dismantled the First Circuit's claim that Maine was merely funding the ``rough equivalent'' of a secular public education. Textually, \S 5204(4) funds tuition at a parent-selected public or private school without requiring private schools to mirror public ones. Operationally, participating private schools can reject applicants, charge tuition well above Maine's cap, employ uncertified teachers, operate single-sex campuses, and design independent curricula if NEASC-accredited. Defining the benefit as ``secular private education'' was a semantic maneuver prohibited by \textit{Agency for Int'l Development v. Alliance for Open Society Int'l, Inc.}, 570 U.S. 205 (2013), that would have gutted \textit{Espinoza}. Nor did the ruling compel Maine to fund religious schooling, since Maine remained free to expand public schools, transportation, tutoring, remote learning, or state boarding schools (\href{https://www.law.cornell.edu/supremecourt/text/20-1088}{Cornell LII}).

\textbf{Rejection of the Status/Use Distinction and Cabining of \textit{Locke v. Davey}:} Turning to the First Circuit's second rationale, the Court emphasized that \textit{Trinity Lutheran} and \textit{Espinoza} never implied that use-based religious discrimination is more permissible than status-based discrimination. Inculcating faith lies at the core of a religious school's mission (\textit{Our Lady of Guadalupe School v. Morrissey-Berru}, 591 U.S. 732 (2020)). Auditing how pervasively a school teaches faith would also generate serious church-state entanglement and denominational favoritism---a line Maine itself barely policed because schools largely self-identified. Finally, the Court restricted \textit{Locke v. Davey} to its historic antiestablishment interest against taxpayer funding of vocational clergy training, finding no comparable historical tradition against aiding private religious schools (\href{https://www.law.cornell.edu/supremecourt/text/20-1088}{Cornell LII}).

\subsection{Dissenting Opinion (Justice Breyer, Joined by Justice Kagan and in Part by Justice Sotomayor)}

Justice Breyer criticized the majority for elevating the Free Exercise Clause while ignoring the Establishment Clause and effectively abandoning the ``play in the joints'' doctrine (\textit{Walz v. Tax Comm'n}, 397 U.S. 664 (1970); \textit{Locke v. Davey}, 540 U.S. 712 (2004)). Because the two Clauses pull in opposite directions, play in the joints affords States legislative leeway to decline to fund religious instruction even when \textit{Zelman} would \textit{permit} such funding. Converting ``may'' into ``must,'' Justice Breyer warned, raised troubling questions about whether public-school districts or charter-school programs must finance religious alternatives (\href{https://www.law.cornell.edu/supremecourt/text/20-1088}{Cornell LII}).

In Part I--B (not joined by Justice Sotomayor), Justice Breyer stressed the Clauses' unifying purpose of preventing religious strife in a pluralistic nation of over 330 million people and more than 100 religious groups, invoking Thomas Jefferson and James Madison against compelled taxpayer support for religious teaching. Distinguishing \textit{Trinity Lutheran} (playground resurfacing) and \textit{Espinoza} (status-based bars on church-controlled schools), he detailed how BCS and Temple Academy weave Christian doctrine throughout science, social studies, and mathematics and apply faith-based admissions and hiring policies. Because Maine's program fulfilled a state constitutional duty to furnish a civic, religiously neutral public education, excluding schools that use public dollars for religious indoctrination fell squarely within \textit{Locke} and avoided inevitable state-versus-school clashes over curricular oversight (\href{https://www.law.cornell.edu/supremecourt/text/20-1088}{Cornell LII}).

\subsection{Dissenting Opinion (Justice Sotomayor)}

Joining all of Justice Breyer's dissent except Part I--B, Justice Sotomayor wrote separately to highlight the speed and sweep of the Court's doctrinal shift since \textit{Trinity Lutheran} in 2017. She argued that the Court had erred from the outset by equating a State's decision not to subsidize religious institutions with unconstitutional discrimination, and noted that her warning in \textit{Trinity Lutheran}---that status-based holdings would soon erase use-based limits---had materialized (\href{https://www.law.cornell.edu/supremecourt/text/20-1088}{Cornell LII}).

Justice Sotomayor emphasized three consequences: (1) discarding the status/use distinction left virtually nothing of the ``play in the joints'' doctrine; (2) Maine taxpayers were now required to finance religious instruction and schools whose policies exclude students and teachers on the basis of religion, sexual orientation, and gender identity; and (3) Maine's tuition program was driven by rural geography rather than an elective preference for private subsidies. While noting that the majority did not appear to compel SAUs to enter \textit{direct} contracts with religious schools (which would exceed \textit{Zelman}'s private-choice rule), she concluded that the Court had transformed separation of church and state from a constitutional commitment into a constitutional violation (\href{https://www.law.cornell.edu/supremecourt/text/20-1088}{Cornell LII}).

\section{Examples: Future Applications}

\begin{enumerate}
  \item \textbf{Same-Side Application (State Education Savings Accounts Excluding Religious Curricula):} A State creates Education Savings Accounts (ESAs) granting parents \$7,000 annually for private-school tuition and educational materials, but bars spending funds at any school that integrates religious teachings into academic courses. Parents wishing to enroll their children in a Catholic high school sue. \textit{Outcome:} \textbf{Unconstitutional under the Free Exercise Clause.} Under \textit{Carson}, disqualifying religious schools from a parent-directed private-school benefit based on religious instruction triggers strict scrutiny and fails.
  \item \textbf{Same-Side Application (Use-Based State Blaine Amendment):} State officials invoke a state constitutional provision prohibiting public money from being ``applied to any religious worship, exercise, or instruction'' to bar Jewish day schools and Islamic academies from a statewide tax-credit scholarship program. \textit{Outcome:} \textbf{Unconstitutional under the Free Exercise Clause.} Because \textit{Carson} eliminated the status/use distinction, use-based Blaine Amendments are just as invalid as status-based ones when applied to exclude religious schools from generally available private-school aid.
  \item \textbf{Opposite-Side Application (Purely Public Rural School Delivery):} A rural State terminates its private-school tuition program and instead constructs regional public high schools, expands busing, opens state-run public boarding dormitories, and offers statewide online public classes. Parents who prefer a private religious academy demand equivalent tuition vouchers. \textit{Outcome:} \textbf{Constitutional.} \textit{Carson} reaffirmed that a State need not subsidize private education at all; the neutrality mandate attaches only when the State chooses to fund private schooling.
  \item \textbf{Opposite-Side Application (Exclusion of Devotional Ministry Degrees):} A State offers merit college scholarships usable at both secular and religious universities, including for academic religion courses, but excludes students pursuing a vocational divinity degree designed to prepare clergy for church leadership. \textit{Outcome:} \textbf{Constitutional.} \textit{Carson} expressly preserved \textit{Locke v. Davey} within its narrow historical domain covering taxpayer funding of vocational ministry training.
  \item \textbf{Fence-Sitter (Generally Applicable Antidiscrimination Conditions or Religious Charter Schools):} After \textit{Carson}, a State removes its ``nonsectarian'' clause from a private-school tuition program but requires all participating private schools---secular and religious alike---to comply with state nondiscrimination laws covering sexual orientation, gender identity, and religion in admissions and hiring. A religious school objects on free-exercise grounds and sues. (A closely related fence-sitter is whether a State must charter a privately operated, publicly funded religious \textit{charter school}, an issue on which the Court later deadlocked 4--4 in \textit{Oklahoma Statewide Charter School Board v. Drummond} (2025).) \textit{Outcome:} \textbf{Uncertain / Fence-sitter.} \textit{Carson} footnote 1* left open whether neutral conduct and antidiscrimination conditions may be enforced against participating religious schools, pitting \textit{Employment Division v. Smith} and civil-rights interests against ministerial-exception autonomy (\textit{Our Lady of Guadalupe}) and \textit{Fulton v. City of Philadelphia} (\href{https://www.cato.org/supreme-court-review/2021-2022/use-nondistinction}{Cato Supreme Court Review}; \href{https://www.acslaw.org/expertforum/carson-v-makin-and-the-dwindling-twilight-of-the-establishment-clause/}{ACS Expert Forum}).
\end{enumerate}

\section{Critique}

\textbf{Scholarly Criticism and Competing Constitutional Premises:} \textit{Carson v. Makin} intensified scholarly debate over how to reconcile the First Amendment's two Religion Clauses. Supporters viewed it as the coherent completion of \textit{Zelman}, \textit{Trinity Lutheran}, and \textit{Espinoza}, whereas critics argued that it subordinated Establishment Clause values to an expansive Free Exercise right to public funding (\href{https://firstamendment.mtsu.edu/article/carson-v-makin/}{First Amendment Encyclopedia}).

\textbf{Progressive Perspective:} Progressive scholars---including Ira C. Lupu and Robert W. Tuttle in the \href{https://www.acslaw.org/expertforum/carson-v-makin-and-the-dwindling-twilight-of-the-establishment-clause/}{ACS Expert Forum} and commentators in the \href{https://houstonlawreview.org/article/92811-unprincipled}{\textit{Houston Law Review}}---contended that \textit{Carson} swept aside 75 years of Establishment Clause precedent forbidding taxpayer financing of worship, proselytizing, and devotional instruction. Lupu and Tuttle noted that Chief Justice Roberts could have ruled narrowly based on \textit{Zelman}'s indirect-aid private-choice framework, but instead discarded the status/use distinction altogether---potentially destabilizing long-settled federal ``Charitable Choice'' rules that bar direct-funded social-welfare contractors (in foster care, adoption, or addiction treatment) from incorporating prayer or worship into publicly funded services. They also stressed an asymmetry problem: if the ministerial exception insulates religious schools from employment-law oversight, the Constitution should not simultaneously compel taxpayers to fund those ecclesiastical activities (\href{https://www.acslaw.org/expertforum/carson-v-makin-and-the-dwindling-twilight-of-the-establishment-clause/}{ACS Expert Forum}). Critics further warned that pressuring States either to finance schools that discriminate against LGBTQ students or to eliminate rural tuition programs harms both public education and religious minorities.

\textbf{Originalist, Textualist, and School-Choice Perspective:} Originalists, textualists, and school-choice proponents---such as petitioners' counsel Michael Bindas in the \href{https://www.cato.org/supreme-court-review/2021-2022/use-nondistinction}{\textit{Cato Supreme Court Review}}---praised \textit{Carson} for honoring the text of the Free Exercise Clause, which protects religious \textit{exercise} rather than merely passive religious identity. Echoing Justice Gorsuch's earlier concurrences, they argued that distinguishing status from use was both unprincipled and entangled state officials in intrusive theological line-drawing that favored secularized religious schools over orthodox ones. Originalist historians also emphasized that state Blaine Amendments stemmed largely from nineteenth-century anti-Catholic nativism rather than founding-era disestablishment principles, when neutral public funding frequently supported religious educational institutions alongside secular ones. Even some conservative federalists, however, noted a lingering structural tension: by shrinking ``play in the joints'' nearly to zero, \textit{Carson} narrows the democratic space for States to adopt stricter church-state funding boundaries in their own constitutions.

\section{Key Quotations}

\begin{quote}
``The State pays tuition for certain students at private schools---so long as the schools are not religious. That is discrimination against religion.''
\end{quote}
\noindent --- Chief Justice Roberts, for the Court (\href{https://www.law.cornell.edu/supremecourt/text/20-1088}{Cornell LII}).

\begin{quote}
``In short, the prohibition on status-based discrimination under the Free Exercise Clause is not a permission to engage in use-based discrimination.''
\end{quote}
\noindent --- Chief Justice Roberts, for the Court (\href{https://www.law.cornell.edu/supremecourt/text/20-1088}{Cornell LII}).

\begin{quote}
``Regardless of how the benefit and restriction are described, the program operates to identify and exclude otherwise eligible schools on the basis of their religious exercise.''
\end{quote}
\noindent --- Chief Justice Roberts, for the Court (\href{https://www.law.cornell.edu/supremecourt/text/20-1088}{Cornell LII}).

\begin{quote}
``The Court today pays almost no attention to the words in the first Clause while giving almost exclusive attention to the words in the second.''
\end{quote}
\noindent --- Justice Breyer, dissenting (\href{https://www.law.cornell.edu/supremecourt/text/20-1088}{Cornell LII}).

\begin{quote}
``Today, the Court leads us to a place where separation of church and state becomes a constitutional violation.''
\end{quote}
\noindent --- Justice Sotomayor, dissenting (\href{https://www.law.cornell.edu/supremecourt/text/20-1088}{Cornell LII}).


\end{document}
